\clearpage
\subsection{LIVE-R1 time contract and broadcast orbit}
\label{sec:live-orbit}

A raw observation is $(t_r^{\rm tag},P_i,\mathrm{PRN}_i,\sigma_{0i})$.
Time is full, unwrapped GPS seconds from 1980-01-06; $P_i$ is the original
GPS L1 C/A pseudorange in metres. Receiver time tags and code share one clock
convention. GPS-labelled RINEX calendar fields receive no UTC conversion.
The final physical reception-time estimate is
\begin{equation}
 \widehat t_r=t_r^{\rm tag}-\widehat b_r/c,
 \qquad c=299\,792\,458\ \mathrm{m\,s^{-1}}.
 \label{eq:live-reception-time}
\end{equation}
Here $b_r$ is the estimated receiver clock bias in metres. The raw tag remains
in the trace. RTCM week rollover is resolved against an explicit reception-date
reference; age tests use full time, never a modulo-week age.

Each GPS ephemeris provides $\sqrt A,e,i_0,\Omega_0,\omega,M_0,\Delta n,
\dot i,\dot\Omega$, six harmonic coefficients, $t_{oe},t_{oc}$, three clock
coefficients, $T_{GD}$, health and issue identifiers. The selected record has
matching PRN, healthy status and $\mathrm{IODC}\bmod256=\mathrm{IODE}$.
Both clock and orbit age are bounded by $7200$ s; a shorter declared half-fit
interval further limits orbit age. A nearer unhealthy issue is not replaced
by an older healthy issue.

For signal emission time $t_s$, define
\begin{align}
 A&=(\sqrt A)^2,&t_k&=t_s-t_{oe},\\
 n&=\sqrt{\mu/A^3}+\Delta n,&M&=M_0+n t_k,
 \qquad \mu=3.986005\times10^{14}\ \mathrm{m^3\,s^{-2}}.
\end{align}
Newton's method solves $E-e\sin E=M$:
\begin{equation}
 E_{j+1}=E_j-\frac{E_j-e\sin E_j-M}{1-e\cos E_j}.
\end{equation}
The implementation reduces $M$ modulo $2\pi$, starts at $E_0=M$, and requires
an increment below $10^{-14}$ rad within 30 iterations. For $0\leq e<1$,
\begin{align}
 \nu&=\operatorname{atan2}\!\left(\sqrt{1-e^2}\sin E,\cos E-e\right),
 &\phi&=\nu+\omega,\\
 u&=\phi+C_{us}\sin2\phi+C_{uc}\cos2\phi,\\
 r&=A(1-e\cos E)+C_{rs}\sin2\phi+C_{rc}\cos2\phi,\\
 i&=i_0+\dot i\,t_k+C_{is}\sin2\phi+C_{ic}\cos2\phi,\\
 \Omega&=\Omega_0+(\dot\Omega-\omega_e)t_k
                  -\omega_e(t_{oe}\bmod604800),\\
 \omega_e&=7.2921151467\times10^{-5}\ \mathrm{rad\,s^{-1}}.
\end{align}
All angles are radians. Radius harmonics are metres; the other harmonics are
angular corrections. These definitions implement the GPS LNAV user algorithm
in \href{https://www.gps.gov/sites/default/files/2025-07/IS-GPS-200N.pdf}{IS-GPS-200N,
section 20.3.3.4.3}; they are new GPS-specific additions to this kernel package.

\clearpage
\subsection{Satellite clock, emission time and reception axes}
\label{sec:live-clock}

Let $x'=r\cos u$ and $y'=r\sin u$. The satellite position in the Earth-fixed
axes at emission is
\begin{equation}
 \widetilde{\boldsymbol x}_s=
 \begin{pmatrix}
 x'\cos\Omega-y'\cos i\sin\Omega\\
 x'\sin\Omega+y'\cos i\cos\Omega\\
 y'\sin i
 \end{pmatrix}.
\end{equation}
With $\Delta t_c=t_s-t_{oc}$, the broadcast clock polynomial, eccentric-orbit
relativity correction, and L1 code clock are
\begin{align}
 \delta t_{\rm poly}&=a_{f0}+a_{f1}\Delta t_c+a_{f2}\Delta t_c^2,\\
 \delta t_{\rm rel}&=F e\sqrt A\sin E,
 &F&=-4.442807633\times10^{-10}\ \mathrm{s\,m^{-1/2}},\\
 \delta t_{s,L1}&=\delta t_{\rm poly}+\delta t_{\rm rel}-T_{GD}.
 \label{eq:live-l1-clock}
\end{align}
The trace stores all three contributions separately. The $-T_{GD}$ sign is
specific to this L1 profile. An external differential C1--P1 bias is not
invented. Alternate code types are rejected until their signal-bias model is
explicit. See \href{https://www.gps.gov/sites/default/files/2025-07/IS-GPS-200N.pdf}{IS-GPS-200N,
sections 20.3.3.3.3.1--20.3.3.3.3.2}.

The measured code connects the receiver and satellite clock tags:
\begin{equation}
 t_s^{\rm tag}=t_r^{\rm tag}-P/c,
 \qquad
 t_s^{(j+1)}=t_s^{\rm tag}-\delta t_{s,L1}(t_s^{(j)}).
 \label{eq:live-transmit-time}
\end{equation}
The adapter uses three clock iterations starting at $t_s^{\rm tag}$. The
receiver clock cancels in $t_r^{\rm tag}-P/c$, because it occurs in both raw
quantities. Subtracting the receiver clock again here would double-correct the
emission time. This construction follows ESA's
\href{https://gssc.esa.int/navipedia/index.php/Emission_Time_Computation}{pseudorange-based
emission-time algorithm}.

Earth rotation over the travel time is applied once, before the native solve:
\begin{align}
 R_3(a)&=\begin{pmatrix}\cos a&\sin a&0\\-\sin a&\cos a&0\\0&0&1\end{pmatrix},\\
 \boldsymbol x_s^{rx}&=R_3(\omega_e\tau)\widetilde{\boldsymbol x}_s,
 &\tau&=\|\boldsymbol x_s^{rx}-\boldsymbol x_r\|/c.
 \label{eq:live-earth-rotation}
\end{align}
Four range iterations, initially $\tau_0=P/c$, determine the reception-frame
satellite position for the current receiver estimate. The native model never
rotates it a second time. The rotation sign and frame convention follow
\href{https://gssc.esa.int/navipedia/index.php/Satellite_Coordinates_Computation}{ESA's
reception-frame coordinate construction}.

The first GSI epoch supplies an independent component check: all eight
emission positions match the native RTKLIB trace within $0.000578$ m and its
clock values within $0.000438$ ns. The RTKLIB trace rounds coordinates to
millimetres and clocks to $0.001$ ns, and excludes $T_{GD}$ from its reported
satellite clock. The comparison therefore adds $T_{GD}$ back to
Equation~\eqref{eq:live-l1-clock} for that check only.

\clearpage
\subsection{Complete Klobuchar L1 ionosphere}
\label{sec:live-ionosphere}

WGS84 ECEF-to-geodetic conversion supplies receiver latitude $\varphi_r$,
longitude $\lambda_r$ and ellipsoidal height $h$. Rotating the line of sight
into local ENU gives azimuth $A_z$ clockwise from north and elevation
$\mathcal E$. A default $10^\circ$ elevation mask is applied after the receiver
geometry has been initialized.

Klobuchar latitude, longitude and elevation arguments are in semicircles:
$\varphi_u=\varphi_r/\pi$, $\lambda_u=\lambda_r/\pi$ and
$\epsilon=\mathcal E/\pi$. Azimuth remains in radians. The complete pierce-point
and geomagnetic construction is
\begin{align}
 \psi&=\frac{0.0137}{\epsilon+0.11}-0.022,\\
 \varphi_I&=\min\!\left(0.416,
       \max\!\left(-0.416,\varphi_u+\psi\cos A_z\right)\right),\\
 \lambda_I&=\lambda_u+\frac{\psi\sin A_z}{\cos(\pi\varphi_I)},\\
 \varphi_m&=\varphi_I+0.064\cos\!\left(\pi(\lambda_I-1.617)\right),\\
 t_{\rm local}&=(43200\lambda_I+t_r^{\rm tag})\bmod86400.
\end{align}
With the four broadcast alpha and four beta coefficients,
\begin{align}
 A_I&=\max\!\left(0,\sum_{j=0}^{3}\alpha_j\varphi_m^j\right),\\
 P_I&=\max\!\left(72000,\sum_{j=0}^{3}\beta_j\varphi_m^j\right),\\
 X&=\frac{2\pi(t_{\rm local}-50400)}{P_I},\\
 F_I&=1+16(0.53-\epsilon)^3.
\end{align}
Amplitude $A_I$ is seconds and period $P_I$ is seconds. The slant L1 delay,
converted to metres, is
\begin{equation}
 I=cF_I
 \begin{cases}
 5\times10^{-9}+A_I\left(1-X^2/2+X^4/24\right),&|X|<1.57,\\
 5\times10^{-9},&\text{otherwise}.
 \end{cases}
 \label{eq:live-klobuchar}
\end{equation}
Every trigonometric operation on a semicircle value explicitly multiplies by
$\pi$. The limits on pierce latitude, amplitude and period are part of the
published model. See
\href{https://gssc.esa.int/navipedia/index.php/Klobuchar_Ionospheric_Model}{ESA's
Klobuchar equations}, and GPS LNAV ionosphere parameters in
\href{https://www.gps.gov/sites/default/files/2025-07/IS-GPS-200N.pdf}{IS-GPS-200N,
section 20.3.3.5.2.5}.

When broadcast coefficients are unavailable, the adapter explicitly selects
$\alpha_j=\beta_j=0$ and records
\code{klobuchar\_zero\_coefficients\_5ns\_baseline}. This still produces
$cF_I\,5\times10^{-9}$ metres of delay; it is not zero ionosphere. For the
internet capture, an independently downloaded BKG broadcast navigation header
provides the coefficients, with its source hash retained. Receiver ephemerides
can still come from the actual RTCM 1019 stream.

\clearpage
\subsection{Troposphere, observation variance and native solve}
\label{sec:live-observation-equation}

LIVE-R1 uses a specified standard atmosphere with relative humidity $H=0.7$.
For $-100\leq h\leq10000$ m and $\mathcal E>0$, define
\begin{align}
 h_0&=\max(h,0),\\
 p&=1013.25(1-2.2557\times10^{-5}h_0)^{5.2568}\quad[\mathrm{hPa}],\\
 T&=288.16-0.0065h_0\quad[\mathrm K],\\
 e_w&=6.108H\exp\!\left(\frac{17.15T-4684}{T-38.45}\right)
       \quad[\mathrm{hPa}].
\end{align}
The hydrostatic and wet Saastamoinen delays are
\begin{align}
 T_h&=\frac{0.0022768p}
 {(1-0.00266\cos2\varphi_r-0.00028h_0/1000)\sin\mathcal E},\\
 T_w&=\frac{0.002277(1255/T+0.05)e_w}{\sin\mathcal E},
 &T_d&=T_h+T_w.
 \label{eq:live-saastamoinen}
\end{align}
Negative terrestrial heights use sea-level weather and carry a specific status.
Outside the height domain, atmosphere is omitted with an explicit diagnostic;
no calibrated weather is inferred. These constants and equations match the
independent \href{https://github.com/tomojitakasu/RTKLIB/blob/71db0ffa/src/rtkcmn.c}{RTKLIB
\code{tropmodel} source}.

The default nominal code uncertainty is $\sigma_{0i}=3$ m. The executed native
profile weights observations by
\begin{equation}
 \sigma_i=\frac{\sigma_{0i}}{\max(\sin\mathcal E_i,0.1)},
 \qquad W_{ii}=\sigma_i^{-2}.
 \label{eq:live-weight}
\end{equation}
CN0 is retained when supplied, without an invented CN0-to-variance calibration.
These are nominal independent variances, not measured position errors or
protection levels.

The full measurement equation and its single additive correction are
\begin{align}
 P_i&=\|\boldsymbol x_r-\boldsymbol x_{s,i}^{rx}\|+b_r
       -c\delta t_{s,i,L1}+I_i+T_{d,i}+\epsilon_i,\\
 a_i&=c\delta t_{s,i,L1}-I_i-T_{d,i},
 &P_{i,\mathrm{corr}}&=P_i+a_i.
 \label{eq:live-corrected-code}
\end{align}
The native CPU/CUDA solver estimates
$\boldsymbol s=(x_r,y_r,z_r,b_r)^\mathsf T$ by weighted least squares. For
$\rho_i=\|\boldsymbol x_r-\boldsymbol x_{s,i}^{rx}\|$,
\begin{equation}
 r_i=P_{i,\mathrm{corr}}-\rho_i-b_r,
 \qquad
 H_i=\begin{pmatrix}
 (\boldsymbol x_r-\boldsymbol x_{s,i}^{rx})^\mathsf T/\rho_i&1
 \end{pmatrix}.
\end{equation}
An online Givens QR solve of the whitened system gives $\Delta\boldsymbol s$;
$\boldsymbol s\leftarrow\boldsymbol s+\Delta\boldsymbol s$. Rank, iteration,
finite-value and residual-fit results remain explicit. Literal ASA/NA selection
is unchanged, including whole-word absorption on a selected boundary hit.

\clearpage
\subsection{Persistent receiver execution and correction feedback}
\label{sec:live-outer-loop}

The persistent native worker receives complete epochs. The preparation layer
and native solve form the outer state-feedback map
\begin{align}
 \mathcal O_n^{(k)}&=\operatorname{Prepare}
       (P_n,t_n^{\rm tag},\mathcal E_n,\boldsymbol s_n^{(k)}),\\
 \boldsymbol s_n^{(k+1)}&=\operatorname{NativeSolve}
       (\mathcal O_n^{(k)},\boldsymbol s_n^{(k)}),
 \label{eq:live-outer-feedback}
\end{align}
where $\mathcal E_n$ denotes the available ephemeris set. Preparation recomputes
reception-frame geometry, elevation selection, ionosphere, troposphere and
weights from the new position. The default outer limit is eight passes.
A valid position requires initialized terrestrial geometry, a converged native
solve, and
\begin{equation}
 \max\!\left(
 \|\boldsymbol x_n^{(k+1)}-\boldsymbol x_n^{(k)}\|,
 |b_n^{(k+1)}-b_n^{(k)}|
 \right)<10^{-3}\ \mathrm m.
\end{equation}

At the first receiver epoch, the complete state is zero. Geometry initialization
uses all available healthy GPS codes with orbit and clock corrections; the
trace explicitly marks atmosphere and elevation as uninitialized. After the
first coarse Earth position is solved, the next pass activates the physical
models. Reference-station coordinates are never read as seeds. Subsequent
epochs start from this worker's preceding valid solution:
\begin{equation}
 \boldsymbol s_{n+1}^{(0)}=\widehat{\boldsymbol s}_n.
\end{equation}
A failed epoch does not overwrite that state. A detected receiver/station change
resets the seed. Duplicate or out-of-order epochs do not produce a new fix.
This is per-epoch code positioning with a warm start; it does not silently add
a motion prior or carrier-phase ambiguities.

The geometric mechanics profile remains a downstream consumer of position
results. Its JK state is not injected into satellite orbit/clock corrections
or used as a fictitious GNSS observation. The live feedback in
Equation~\eqref{eq:live-outer-feedback} closes the formerly missing receiver
preparation loop around the existing position kernel.

\begin{panel}
\textbf{Separate observable evidence.}
The journal stores raw measurements, every preparation pass, every native
solution and final availability. Correction components retain signs and units.
A recorded-file replay, a local timed replay, and a current internet receiver
stream are distinct source modes. A remote reference receiver's position is
its antenna position; it does not measure the position of the user's computer.
\end{panel}

The GSI 120-epoch run starts at ECEF zero, converges on all epochs, and agrees
with independent NumPy solves of all 263 converged passes within
$1.17\times10^{-8}$ m per state component. The executed weighting differs from
RTKLIB's broadcast-variance weighting: their final positions differ by a mean
$0.529$ m and a maximum $1.589$ m. A controlled NumPy re-solve with the same
physical corrections and RTKLIB weights reduces those differences to a mean
$0.000174$ m and a maximum $0.000346$ m. This experiment isolates the weighting
choice; it is recorded separately from the executed native profile. The full
report is \path{results/real_0759_cpu/independent_comparison.json}.
